\documentclass[10pt,aspectratio=169]{beamer}

% --- Minimal theme (matches 01_distance_slides.tex) ------------------------
\usetheme{default}
\usecolortheme{dove}
\usefonttheme{professionalfonts}
\setbeamertemplate{navigation symbols}{}
\setbeamertemplate{footline}[frame number]
\setbeamertemplate{itemize items}[circle]
\setbeamertemplate{section in toc}[sections numbered]
\setbeamerfont{frametitle}{size=\large,series=\bfseries}
\setbeamercolor{frametitle}{fg=black}
\setbeamercolor{title}{fg=black}
\setbeamercolor{structure}{fg=black}
\setbeamercolor{alerted text}{fg=black!80}

\usepackage{booktabs}
\usepackage{array}
\usepackage{graphicx}
\usepackage{xcolor}
\graphicspath{{../figs/}}

\definecolor{purpleC}{HTML}{762A83}
\definecolor{greenC}{HTML}{1B7837}

\title{Country-by-country attribution of cross-reanalysis disagreement}
\subtitle{Generalising deck 01 to all six reanalysis pairs (BHM / GDP-pc)}
\author{}
\date{\today}

\begin{document}

\begin{frame}\titlepage\end{frame}

\begin{frame}{Outline}\tableofcontents\end{frame}

% ===========================================================================
\section{Setup and metric}

\begin{frame}{What we're attributing}
  \begin{itemize}
    \item Same BHM growth panel, same specification, four reanalyses
          (\textbf{ERA5, MERRA2, JRA-3Q, GMFD}).
    \item Each delivers a marginal-effect curve
          $\mathrm{ME}(T) = \beta_1 + 2\beta_2 T$.
    \item \textbf{Question}: which countries drive the cross-reanalysis
          disagreement in the fitted dose-response?
    \item \textbf{Deck 01} answered this for the single \emph{pair}
          ERA5-vs-MERRA2. This deck generalises to the mean across all
          \emph{six} pairs and looks per pair separately.
  \end{itemize}
\end{frame}

\begin{frame}{The L1 distance, and where the numbers come from}
  For any two reanalyses $a, b$:
  \[
    L1(a,b) \;=\; \int \bigl| \mathrm{ME}_a(T) - \mathrm{ME}_b(T)\bigr|\, f(T)\,dT
  \]
  where $f(T)$ is the empirical \textbf{population-weighted} distribution of
  annual mean T over the regression sample.

  \vspace{0.6em}
  \textbf{Computed by Kevin}
  (\texttt{code/3b\_dmgfdiffs\_sources.ipynb}, function \texttt{calc\_dmes}).
  Outputs:
  \begin{itemize}\setlength{\itemsep}{1pt}
    \item \texttt{dmes/dmes\_gdppc.csv} --- per (reanalysis pair, dropped
          country); 6 pairs $\times$ 196 countries $= 1{,}176$ rows.
    \item \texttt{dmes/mean\_dmes\_gdppc.csv} --- mean across pairs per LOO.
    \item Each file contains a special row \texttt{loo\_id == "all"}
          giving the full-sample baseline. Sign of $\%$ reduction in this
          deck:
          $\text{pct} = (L1_\text{full} - L1_\text{LOO}) / L1_\text{full}$
          $\times 100$.
          \textcolor{purpleC}{Positive $\Rightarrow$ drives disagreement};
          \textcolor{greenC}{negative $\Rightarrow$ drives agreement}.
  \end{itemize}
\end{frame}

% ===========================================================================
\section{Walkthrough: how the L1 distance is computed (Bahrain example)}

\begin{frame}{Step 0 -- marginal-effect curves with and without Bahrain}
  \centering
  \includegraphics[height=0.8\textheight]{bahrain_me_curves.pdf}\\[-0.4em]
  \footnotesize Each panel: BHM dose-response fit on the full sample (solid)
  vs.\ refit after dropping every Bahrain observation (dashed). Coefficients
  are from \texttt{reg\_outputs/gdppc/regoutputs\_coeffs\_*\_no\{Bahrain\}.csv}
  (computed by Kevin).
\end{frame}

\begin{frame}{Contrast: same view but for Greenland (a ``drives agreement'' country)}
  \centering
  \includegraphics[height=0.78\textheight]{greenland_me_curves.pdf}\\[-0.4em]
  \footnotesize Bahrain pulled the four MEs together; Greenland is the
  opposite case --- removing it makes ERA5 / JRA-3Q / GMFD diverge from
  MERRA2 (which barely moves). That's why Greenland shows up green
  (drives agreement) on the averaged map.
\end{frame}

\begin{frame}{Walkthrough: full sample}
  \centering
  \includegraphics[height=0.82\textheight]{bahrain_l1_walkthrough_full.pdf}
\end{frame}

\begin{frame}{Walkthrough: after Bahrain is dropped}
  \centering
  \includegraphics[height=0.82\textheight]{bahrain_l1_walkthrough_noBahrain.pdf}\\[-0.4em]
  \footnotesize The shaded purple area (the integrand
  $|\Delta\mathrm{ME}| \cdot f(T)$) shrinks slightly.
  Computed here in my own code (\texttt{tb/code/bahrain\_walkthrough.py}) using
  Kevin's coefficients; the resulting full-sample L1 = 0.00535 matches
  Kevin's CSV baseline (0.00537) within rounding.
\end{frame}

% ===========================================================================
\section{Per-pair country attribution}

\begin{frame}{Per-pair LOO maps + histograms (the main new view)}
  \centering
  \includegraphics[height=0.85\textheight]{loo_perpair_maps.pdf}
\end{frame}

\begin{frame}{Same view, drivers only (negative countries collapsed to grey)}
  \centering
  \includegraphics[height=0.85\textheight]{loo_perpair_maps_positive.pdf}\\[-0.3em]
  \footnotesize Only countries with positive pct (i.e.\ dropping them
  reduces the pair's L1) are coloured. The histograms below show what
  was hidden: purple bars are the positive countries that survive on
  the map; grey bars are the negative countries that get folded into the
  grey background.
\end{frame}

\begin{frame}{Reading the per-pair maps}
  \begin{itemize}\setlength{\itemsep}{2pt}
    \item Each of the 6 maps is a separate reanalysis pair. Each map is on
          its \emph{own} colour scale because the per-pair baseline L1
          varies $\sim$4.6$\times$ (MERRA2-GMFD = 0.0015; JRA-3Q-GMFD = 0.0069).
    \item Histogram below each map: distribution of the same percentage
          across all 195 countries for that pair.
    \item Annotations name the top-3 drivers and bottom-3 ``glue''
          countries per pair.
    \item \textbf{Deck 01 finding} (ERA5 vs MERRA2, mid panel): Sierra Leone,
          Papua New Guinea, Chad are the top drivers --- still visible here.
    \item \textbf{New finding} (other pairs): Bahrain, Guinea-Bissau,
          Marshall Islands, Micronesia appear at the top of pairs that
          deck 01 never considered (GMFD-vs-ERA5, JRA-3Q-vs-GMFD).
  \end{itemize}
\end{frame}

% ===========================================================================
\section{Where in the world does the disagreement live?}

\begin{frame}{Scatter: country influence vs.\ mean T and mean income}
  \centering
  \includegraphics[width=0.95\linewidth]{loo_scatter.pdf}\\[-0.2em]
  \footnotesize
  Disagreement-driving countries (purple) cluster in the hot tail
  (\textgreater 25\textdegree C) and at the lower-to-middle end of the income
  distribution. ``Glue'' countries (green) span the whole income range but
  Greenland and Liberia stand out as extreme outliers.
\end{frame}

\begin{frame}{Per-pair scatter vs.\ mean temperature}
  \centering
  \includegraphics[height=0.78\textheight]{loo_scatter_perpair_T.pdf}\\[-0.3em]
  \footnotesize One panel per reanalysis pair, $y$-axis = \% reduction in
  \emph{that pair's} L1 when each country is dropped. Each panel uses its
  own colour norm centred at 0. The MERRA2-vs-GMFD panel (top-left) has
  the largest swings because that pair's baseline L1 is smallest --- small
  absolute coefficient shifts go a long way percentage-wise.
\end{frame}

\begin{frame}{Per-pair scatter vs.\ mean log income}
  \centering
  \includegraphics[height=0.78\textheight]{loo_scatter_perpair_income.pdf}\\[-0.3em]
  \footnotesize Same layout, $x$-axis = country-mean log GDP-pc PPP.
  Low-income tropical countries (Liberia, Sierra Leone, Marshall Islands,
  Papua New Guinea) recur as drivers across multiple pairs; high-income
  cold-region outliers (Greenland, UAE) tend to glue specific pairs
  together.
\end{frame}

\begin{frame}{Per-pair residual-T disagreement vs.\ climate and income}
  \centering
  \includegraphics[width=\linewidth]{tc_disagr_perpair.pdf}\\[-0.3em]
  \footnotesize
  \emph{Different metric, same axes:} $y$ = pairwise MAD of residualised
  temperature time series (resid from
  \texttt{T $\sim$ 1 | iso + year + iso[t] + iso[t$^2$]} per reanalysis,
  matching the deck-02 R code).
  Unlike the LOO ME-coefficient influence (previous slides, $R^2 \approx 0$),
  the underlying temperature disagreement \emph{does} sort by income and
  climate: low-income, hot countries see systematically larger
  inter-reanalysis residual disagreement. The message: measurement noise is
  income-stratified; the part of it that propagates to fitted
  $\beta_1, \beta_2$ is dominated by a handful of tail-T leverage points.
\end{frame}

% ===========================================================================
\section{Raw and residual temperature series for the panel countries}

\begin{frame}{Raw annual T -- top 3 disagreement, bottom 3 agreement}
  \centering
  \includegraphics[height=0.82\textheight]{country_temp_raw.pdf}\\[-0.4em]
  \footnotesize Population-weighted annual mean T per reanalysis. Each panel
  has its own y-axis (T levels vary $\sim$30\textdegree C across the six
  countries).
\end{frame}

\begin{frame}{Residualised annual T (same panels)}
  \centering
  \includegraphics[height=0.82\textheight]{country_temp_residual.pdf}\\[-0.4em]
  \footnotesize $T_{\mathrm{resid}} = T - \widehat T$ where
  $\widehat T$ is the in-sample fit of
  $T \sim \text{country FE} + \text{year FE} + \text{country-specific linear} +$
  $\text{country-specific quadratic time trends}$
  (Tom's R spec in \texttt{tb/code/explore/temp\_compare.R};
  re-implemented in Python with sparse OLS in
  \texttt{tb/code/country\_residuals.py}). This is the variation the BHM
  regression actually identifies $\beta_1, \beta_2$ from.
\end{frame}

\begin{frame}{What the residual panels show}
  \begin{itemize}\setlength{\itemsep}{2pt}
    \item \textbf{Bahrain, Guinea-Bissau, Papua New Guinea} (drivers): the
          four reanalyses already disagree at the year-to-year level
          \emph{after} all trends are removed. Year-on-year noise differs
          by reanalysis, which is what the BHM coefficients pick up.
    \item \textbf{Greenland, Liberia} (glue): also have residual
          disagreement (GMFD diverges sharply in Liberia in the late 2000s
          -- the same anomaly Tom flagged in deck 02). But the structure
          of that disagreement happens to \emph{average out} across pairs in
          a way that keeps the dose-responses aligned.
    \item Takeaway: who ``drives'' vs ``glues'' the disagreement is
          determined jointly by where the residual noise is large
          \emph{and} where it covaries with growth.
  \end{itemize}
\end{frame}

% ===========================================================================
\section{Sensitivity: pop-weighted vs uniform integration}

\begin{frame}{Why we might worry about pop weighting}
  \begin{itemize}
    \item Kevin's headline L1 integrates $|\Delta\mathrm{ME}|$ against the
          \emph{pop-weighted} empirical T density. India / China dominate
          the integral.
    \item An unweighted (each $(country, year, reanalysis)$ row equal)
          integral asks a more abstract regression-stats question:
          ``which countries are influential in the OLS fit?''
    \item Calibration: pop recompute (my code) vs Kevin's CSV →
          Spearman 0.92, mean ratio 0.92 (small offset because Kevin's
          xarray uses the BCD T grid; my code uses the BHM panel).
          Both pop and uniform paths come from the shared
          \texttt{tb/code/funcs\_dme.py} helper.
  \end{itemize}
\end{frame}

\begin{frame}{Main map: pop-weighted vs uniform}
  \begin{columns}[T]
    \begin{column}{0.5\textwidth}
      \centering pop-weighted (Kevin's CSV)\\[-0.1em]
      \includegraphics[width=\linewidth]{loo_country_influence.pdf}
    \end{column}
    \begin{column}{0.5\textwidth}
      \centering uniform (each obs equal)\\[-0.1em]
      \includegraphics[width=\linewidth]{loo_country_influence_uniform.pdf}
    \end{column}
  \end{columns}\vspace{0.2em}
  \footnotesize Top-3 unchanged (Bahrain / Guinea-Bissau / PNG).
  Bottom-3 swap Trinidad and Tobago for UAE. Range tightens from
  $[-14.6\%, +5.5\%]$ pop → $[-10.8\%, +5.6\%]$ uniform.
\end{frame}

\begin{frame}{Per-pair maps: pop vs uniform (positive-only variant)}
  \begin{columns}[T]
    \begin{column}{0.5\textwidth}
      \centering pop-weighted\\[-0.1em]
      \includegraphics[width=\linewidth]{loo_perpair_maps_positive.pdf}
    \end{column}
    \begin{column}{0.5\textwidth}
      \centering uniform\\[-0.1em]
      \includegraphics[width=\linewidth]{loo_perpair_maps_positive_uniform.pdf}
    \end{column}
  \end{columns}\vspace{0.2em}
  \footnotesize Same six pairs, same pair ordering. The MERRA2-vs-GMFD
  panel (top-left) is where weighting matters most: that pair has the
  smallest baseline so small absolute coefficient shifts translate to
  large \%-swings regardless of weighting.
\end{frame}

\begin{frame}{Same view but \emph{absolute} reduction (shared L1 scale across pairs)}
  \begin{columns}[T]
    \begin{column}{0.5\textwidth}
      \centering pop-weighted\\[-0.1em]
      \includegraphics[width=\linewidth]{loo_perpair_maps_positive_abs.pdf}
    \end{column}
    \begin{column}{0.5\textwidth}
      \centering uniform\\[-0.1em]
      \includegraphics[width=\linewidth]{loo_perpair_maps_positive_abs_uniform.pdf}
    \end{column}
  \end{columns}\vspace{0.2em}
  \footnotesize
  Now $y$ = \emph{absolute} reduction in pair-L1 when country dropped (same
  L1 units across all 6 pairs $\Rightarrow$ one shared colour scale and one
  shared histogram x-axis). The high-$L^1$-baseline pairs (MERRA2-vs-JRA-3Q,
  JRA-3Q-vs-GMFD, GMFD-vs-ERA5) carry the biggest absolute movements; the
  low-baseline MERRA2-vs-GMFD panel that looked dramatic in the \% view
  shrinks to almost nothing. Inverts the takeaway about ``where the
  swings are biggest.''
\end{frame}

\begin{frame}{Averaged scatter: pop vs uniform}
  \begin{columns}[T]
    \begin{column}{0.5\textwidth}
      \centering pop-weighted\\[-0.1em]
      \includegraphics[width=\linewidth]{loo_scatter.pdf}
    \end{column}
    \begin{column}{0.5\textwidth}
      \centering uniform\\[-0.1em]
      \includegraphics[width=\linewidth]{loo_scatter_uniform.pdf}
    \end{column}
  \end{columns}\vspace{0.2em}
  \footnotesize Cluster of disagreement-drivers in the hot tail survives
  intact under uniform weighting. The income relationship is also
  qualitatively unchanged.
\end{frame}

\begin{frame}{Sensitivity takeaway}
  \begin{itemize}\setlength{\itemsep}{4pt}
    \item Spearman rank correlation (uniform vs Kevin pop) = \textbf{0.99}.
          The country ranking is essentially robust to the weighting choice.
    \item ``Tropical hot-tail drives disagreement'' is \emph{not} a
          pop-weighting artefact -- it survives uniform integration.
    \item What does shift: Sierra Leone drops out of the top-5 under
          uniform (replaced by Equatorial Guinea); Greenland's negative
          magnitude shrinks from $-14.6\%$ to $-10.8\%$ (pop weighting
          gave Greenland near-zero density, so dropping it barely moved
          the integrand; uniform gives Greenland \textgreater 100$\times$
          more weight).
    \item All variants are produced by a single \texttt{funcs\_dme.py}
          helper; flip \texttt{weighting="uniform"}/\texttt{"pop"} to
          regenerate any figure.
  \end{itemize}
\end{frame}

% ===========================================================================
\section{Takeaways}

\begin{frame}{Takeaways}
  \begin{itemize}\setlength{\itemsep}{4pt}
    \item Per-pair attribution generalises deck 01 cleanly. Sierra Leone,
          Papua New Guinea, Chad remain top drivers \emph{within the
          ERA5-vs-MERRA2 pair}; Bahrain, Guinea-Bissau emerge as top drivers
          for pairs involving GMFD and JRA-3Q.
    \item The all-pairs average highlights Bahrain, Guinea-Bissau, PNG --
          countries that drive disagreement in \emph{multiple} pairs.
          Sierra Leone drops from \#1 (pairwise) to \#5 (averaged).
    \item Disagreement concentrates in (hot, low-to-mid income) countries.
          Greenland and Liberia are large negative outliers but in different
          income brackets.
    \item Residualised T panels show the year-to-year reanalysis disagreement
          that drives the dose-response disagreement. GMFD's late-2000s
          divergence (visible in Liberia, Guinea-Bissau) is the dominant
          source of GMFD-vs-others differences.
  \end{itemize}
\end{frame}

\begin{frame}{What is Kevin's, what is mine}
  \begin{itemize}\setlength{\itemsep}{3pt}
    \item \textbf{Kevin}: regression coefficients (full and LOO) in
          \texttt{reg\_outputs/gdppc/}; per-pair and averaged dme tables in
          \texttt{dmes/}.
    \item \textbf{Mine}: per-pair maps + histograms; Bahrain walkthrough
          (re-computed L1, matches Kevin's baseline within rounding);
          scatter plots; Python re-implementation of Tom's R residualisation.
    \item All scripts in \texttt{tb/code/loo\_*.py},
          \texttt{tb/code/bahrain\_walkthrough.py},
          \texttt{tb/code/country\_residuals.py}.
    \item Slide deck source: \texttt{tb/slides/06\_loo\_country\_attribution\_slides.tex}.
  \end{itemize}
\end{frame}

\end{document}
